\section{Counting the sections carried by a component}
\label{sec:rank-share}
Let \(Z\) be a component at the chosen deep level, \(D\) the dimension of
its flag variety, and \(T\) the integer in the weight construction. Let
\(\delta>0\) be the degree of the image of the cusp linear system of
weight \(u\) on \(Z\). The source uses the birationality of this system.
The cusp dimension at weight \(\lambda\) is \(d_\lambda\). We prove
the estimate used in the height calculation:
\begin{equation}\label{eq:rank-share-count}
 \left(1-7(D^2/T+D\rho)\right)d_\lambda
 \ \le\ \binom{T-1}{D}\delta .
\end{equation}
For the terminal weight \(\lambda_T\) there is also
\(\binom{T-1}{D}\delta\le d_{\lambda_T}\).
The upper inequality applies to all the weights satisfying the dominance
conditions below. The lower one uses sections of a particular adjoint bundle.

\paragraph{Geometric data.}
Assume \(p,q\ge1\) and \([E:\mathbb Q]/2>2\), as in the selection of the
component. Here \(D\ge2\), \(T\ge D+1\), and \(0\le\rho\le1/2\).
Write \(v\) for the top-degree cusp growth of \(u\). The proportionality
and valuation-multiplicity comparison gives
\begin{equation}\label{eq:rankshare-vol}
 (1-\rho)v\le\delta .
\end{equation}
This comparison uses the strictly decreasing weights \(u,u'\) and the
equality of component volumes at deep level. Put \(n=T+D+3\). At each
embedding assume that both \(\lambda\) and \(nu-\lambda\) are weakly
decreasing, and that all consecutive gaps of \(u\) are at least
\[
 (D+3)\left(\binom{D+1}{2}+4(p+q)+D+3\right).
\]
The geometric inputs in the source give the following model.
On a projective modification \(X\) of the compactified flag component the
linear system of \(u\) has generated moving bundle \(H\) and effective
fixed divisor \(\Delta\). Write \(U=H+\Delta\) and \(B=L_\lambda\)
in additive line-bundle notation, using pullbacks on \(X\). Both
\(U,B\) and also \(nU-B\) are nef. The morphism of \(H\) is birational
onto the component image, so the projection formula gives \(H^D=\delta\).
The comparison between cusp growth and nef intersection gives \(U^D\le v\).
Cusp forms extend across the compactification and inject into \(H^0(X,B)\).
These assertions about compactification, growth and nefness are the
geometric hypotheses of the following calculation.

\paragraph{Upper bound by hyperplane sections.}
For \(a\ge0\), an integral projective \(D\)-fold, a generated \(H\) of
top intersection \(\delta>0\), and a line bundle \(B\), we claim
\[
 B H^{D-1} < (a+1)\delta
 \quad\Longrightarrow\quad h^0(B)\le\binom{a+D}{D}\delta .
\]
Take a generic member \(X_1\) of the linear system after extending scalars
to the rational function field of its coefficients. Cohomology dimensions
are preserved by this field extension. The generic divisor is integral
over that field when nonempty: on each open where one generating section
is a unit, the universal linear equation eliminates its coefficient.
The incidence scheme is locally an affine bundle over the integral base,
and its generic fiber is a localization. A nonempty section drops the
dimension by one. Restriction to it preserves \(\delta\) as
\((H|_{X_1})^{D-1}\). For dimension one the restriction has length
\(\delta\) over the coefficient field. The exact sequence for the section
bounds \(h^0(B)\) by the sum of \(h^0(B-H)\) and
\(h^0(X_1,B|_{X_1})\).

For clarity, in dimension \(r>0\) use the parameter \(b=a+1\). If
\(BH^{r-1}<0\), a nonzero section would have an effective zero divisor
whose intersection with \(r-1\) nef classes is nonnegative. Thus the
bound for \(b=0\) is zero. Passing to \(B-H\) subtracts \(\delta\)
from the mixed intersection and replaces \(b\) by \(b-1\). Restriction
to \(X_1\) leaves \(b\) fixed and replaces \(r\) by \(r-1\), since
\((B|_{X_1})(H|_{X_1})^{r-2}=BH^{r-1}\) for \(r\ge2\).
In dimension zero the section dimension of an invertible sheaf equals
the length, giving \(\delta\). For \(b\ge1\), induction in \(r,b\)
therefore gives
\[
 \left(\binom{b+r-2}{r}+\binom{b+r-2}{r-1}\right)\delta
   =\binom{b+r-1}{r}\delta
\]
whenever \(b+r\ge2\), using the zero initial value for \(b=0\).
This proves the claim.

On the modification, replacing \(H\) by \(U\) one factor at a time
in a product with other nef factors can only increase the intersection,
because \(U-H=\Delta\) is effective. Nefness of \(nU-B\) gives
\[
 BH^{D-1}\le nUH^{D-1}\le nU^D\le nv .
\]
Take \(a=\lfloor n/(1-\rho)\rfloor\). Inequality
\eqref{eq:rankshare-vol} and \(\delta>0\) imply
\(nv<(a+1)\delta\), including when \(n/(1-\rho)\) is an integer.
Thus \(d_\lambda\le\binom{a+D}{D}\delta\).

\paragraph{The loss in replacing \(a+D\) by \(T-1\).}
Write \(N=a+D\), \(M=T-1\), and \(x=4D/T+\rho\).
We have \(D\le M\le N\), \(a\ge n\), and
\(a(1-\rho)\le T+D+3\). Consequently
\[
 N-M\le \rho a+2D+4\le x(a+1)=x(N-D+1).
\]
Indeed \(2D+4\le4D\) and \(4D(a+1)/T\ge4D\).
If \(x\le1\), then for \(0\le j<D\)
\((M-j)/(N-j)\ge1-x\); multiplying gives
\(\binom{M}{D}\ge(1-Dx)\binom{N}{D}\) by Bernoulli's inequality.
The last inequality also holds when \(x>1\), since its right side is
nonpositive. Finally \(Dx=4D^2/T+D\rho\le7(D^2/T+D\rho)\).
If the coefficient on the left of \eqref{eq:rank-share-count} is
nonpositive the result is immediate. Otherwise multiply the upper bound
for \(d_\lambda\) by that coefficient and apply the binomial estimate.

\paragraph{Lower bound at the terminal weight.}
In the component data, the generated moving line bundle \(H\) has
\(H^D=\delta\). The adjoint sections of
\(K_X+A+TH\) inject into the terminal cusp space under the chosen frame.
The auxiliary divisor \(A\) is nef and big. We use the source's
Kawamata--Viehweg vanishing for
\(K_X+A+jH\), \(j\ge1\), in positive cohomological degrees on the
smooth projective model in characteristic zero. All identifications of
these ordinary sections with subspaces of cusp forms are made before
counting them.

Here is the counting lemma. Let a coherent sheaf \(F\) have support
dimension at most \(r\), and suppose that every positive cohomology group
of \(F(jH)\) vanishes for \(j\ge j_0\). For \(m\ge0\),
\[
 h^0(F((j_0+m)H))\ge \binom{m}{r}(H^r\cdot F).
\]
In dimension zero the dimension of sections equals the length over the
ground field and twisting preserves it. For \(r>0\), choose a section of
\(H\) outside the finitely many associated primes of \(F\). This is
possible over the infinite coefficient field because \(H\) is generated.
Multiplication by it is injective. The quotient \(F_1\) after restriction
has support dimension at most \(r-1\), with
\(H^{r-1}\cdot F_1=H^r\cdot F\).
The long exact cohomology sequence gives the same vanishing for \(F_1\)
starting at \(j_0+1\), and gives
\[
 h^0(F((j+1)H))-h^0(F(jH))=h^0(F_1((j+1)H)),\quad j\ge j_0 .
\]
Summing the dimension-\((r-1)\) bound for \(m\) successive differences,
and using \(\sum_{i=0}^{m-1}\binom{i}{r-1}=\binom{m}{r}\), proves
the claim. Apply it to \(F=K_X+A,\ j_0=1,\ m=T-1,\ r=D\). Twisting
the structure sheaf by an invertible sheaf preserves the top intersection.
The section injection therefore gives \(d_{\lambda_T}\ge
\binom{T-1}{D}\delta\), as required.

\paragraph{Source correspondences.}
In the original \srcid{DivisorProofs} directory the geometric reductions are
\srcid{RankComponentWeightInvariance} and
\srcid{ShareStatementOfKnownTheorems}. The binomial comparison is
\srcid{ChooseMulOneSubLeChoose}. For the lower bound see
\srcid{GlobalSectionRankLowerBoundOfHilbertSamuel},
\srcid{ChooseMulIntersectionNumberLeSectionDim}, and
\srcid{ChooseMulComponentDegreeLeDimension}.
